\chapter{Sistemas de Detección y Alarma de Incendios}
\label{sec.cap5:sistemas_deteccion}

El tiempo transcurrido entre la ignición incipiente y la detección efectiva determina el éxito o fracaso de cualquier
estrategia de protección contra incendios. Los sistemas automáticos de detección y alarma, normalizados bajo los
lineamientos de las normas NFPA 72 e IRAM 3554, tienen por finalidad identificar el siniestro en su fase embrionaria,
alertar a los ocupantes para permitir su evacuación antes de que se alcancen concentraciones letales de gases tóxicos, y
comandar las acciones automáticas de supresión y compartimentación.

\section{Arquitectura y Topología de los Sistemas de Detección}

Un sistema integral de detección se compone de paneles de control, sensores
automáticos de campo, pulsadores manuales de alarma y dispositivos de notificación acústico-óptica. Según la topología de
comunicación y la granularidad de la información transmitida, se distinguen dos arquitecturas principales:

\begin{itemize}
    \item Sistemas Convencionales: Agrupan los detectores y avisadores en circuitos o zonas físicas cableadas en
          configuración radial con resistencia de fin de línea. Cuando un detector conmuta a estado de alarma, la central
          detecta un incremento de corriente provocado por la baja impedancia del dispositivo activado. Este esquema
          presenta un costo inicial reducido, pero posee la severa limitación operativa de no individualizar el sensor en
          falla o alarma dentro de la zona, requiriendo una búsqueda visual en todo el sector comprometido. Su aplicación se
          restringe a locales de reducida superficie o naves de ambiente único.
    \item Sistemas Analógicos-Direccionables: Emplean lazos de comunicación digital bidireccional en configuración de
          anillo cerrado (Lazos Clase A) o abiertos (Clase B). Cada dispositivo de campo dispone de una dirección binaria
          propia asignada por microinterruptores o programación electrónica. Los sensores no operan como simples
          conmutadores de umbral, sino que reportan periódicamente a la central un valor analógico continuo proporcional a
          la magnitud física monitoreada (porcentaje de oscurecimiento por metro, temperatura puntual, concentración de
          partículas). Esta tecnología permite que el panel ejecute algoritmos de compensación automática por suciedad
          acumulada, genere prealarmas programables y localice con exactitud milimétrica el punto del siniestro,
          reduciendo drásticamente las falsas alarmas y facilitando el mantenimiento predictivo.
\end{itemize}

\section{Detectores Ópticos de Humo (Fotoeléctricos)}

Los detectores ópticos puntuales constituyen la tecnología más difundida en instalaciones comerciales y corporativas. Su
principio de funcionamiento se sustenta en una cámara de muestreo laberíntica oscura que impide el ingreso de la luz
ambiental pero permite la libre circulación del aire. En el interior de la cámara se ubica un diodo emisor de luz
infrarroja (LED) y un fotodiodo receptor colocado en ángulo diedro (típicamente entre $90^\circ$ y $120^\circ$) respecto
al eje óptico de emisión.

En condiciones de aire limpio, el haz luminoso se propaga en línea recta sin incidir sobre el fotodiodo receptor,
manteniendo la corriente de reposo en valores prácticamente nulos. Cuando las partículas de humo ingresan al laberinto,
interactúan con el haz luminoso produciendo el fenómeno de dispersión óptica de la luz (efecto Tyndall). Una porción de
los fotones dispersados impacta sobre el fotodiodo sensible, generando una corriente proporcional a la densidad óptica
del aerosol.

Esta tecnología resulta altamente idónea para detectar incendios de combustión lenta, incompleta o latente, los cuales
generan aerosoles con partículas visibles de diámetro medio a grande (entre $0.3\,\mu\mathrm{m}$ y $10\,\mu\mathrm{m}$),
tales como los generados por el sobrecalentamiento de cables con aislación de PVC o la pirolización de espumas y
celulosas. Sus limitaciones radican en la vulnerabilidad al vapor de agua, condensaciones y acumulación de polvillo
industrial, requiriendo mantenimiento periódico de limpieza de cámara.

\section{Detectores Iónicos de Humo}

Los detectores iónicos basan su principio físico en la conducción eléctrica de gases ionizados. Disponen de una cámara de
muestreo que aloja una diminuta fuente radiactiva encapsulada de Americio-241 ($^{241}\mathrm{Am}$), con una actividad
del orden de $1\,\mu\mathrm{Ci}$, emisora de partículas alfa ($\alpha$). Estas partículas ionizan las moléculas neutras de
nitrógeno y oxígeno del aire circundante, escindiéndolas en pares de iones positivos y electrones libres.

Bajo la acción de una diferencia de potencial aplicada entre dos electrodos metálicos polarizados, los iones migran hacia
las placas opuestas, estableciendo una corriente iónica continua de reposo sumamente pequeña, en el rango de los
picoamperes ($p\A$). Cuando partículas de combustión submicrónicas invisibles (con diámetros inferiores a
$0.3\,\mu\mathrm{m}$) penetran en la cámara, los iones se adhieren a la masa de las partículas carbonosas debido a su mayor
inercia y menor movilidad electroforética. Esto reduce la velocidad de recombinación y disminuye drásticamente la
corriente eléctrica de la cámara. Al caer la corriente por debajo del umbral de disparo, la electrónica activa la
señal de alarma.

Los detectores iónicos poseen una velocidad de respuesta superior ante fuegos rápidos con llama abierta y escaso humo
visible. No obstante, en la actualidad su uso se encuentra prácticamente descontinuado a escala global debido a las
restricciones ambientales y regulatorias asociadas a la gestión, transporte y disposición final de fuentes radiactivas,
siendo sustituidos por sensores ópticos avanzados de longitud de onda dual.

\section{Detectores Lineales de Humo por Haz Proyectado (Barreras Ópticas)}

En naves industriales de gran altura, centros logísticos, plantas de producción y hangares donde los techos superan los
$6\m$ a $12\m$, los detectores puntuales de techo pierden eficacia debido al fenómeno de estratificación térmica. La masa
ascendente de humo y gases calientes pierde flotabilidad a medida que se mezcla con el aire ambiental más frío,
formando una capa horizontal de estratificación suspendida que no alcanza la cumbrera del techo.

Para resolver este desafío de ingeniería se implementan los detectores lineales de haz proyectado . El sistema está
integrado por un emisor de radiación infrarroja y un receptor fotosensible colocados en muros opuestos a distancias
operativas de entre $10\m$ y $100\m$, o bien un transceptor monobloque que emite y recibe el haz reflejado desde un
prisma catadióptrico retrorreflector.

El principio de detección es el oscurecimiento óptico a lo largo de una trayectoria extendida: el sensor mide de manera
continua la atenuación del flujo luminoso recibido respecto al calibrado en reposo. Cuando la capa de humo estratificado
cruza la trayectoria del haz, el porcentaje de atenuación sobrepasa el umbral programado (e.g., entre $20\percent$ y
$50\percent$ de absorción) durante un tiempo de integración fijado, disparando la condición de incendio. Asimismo, estos
equipos disponen de algoritmos de discriminación que identifican el bloqueo súbito y total del haz (provocado por una grúa
puente o carretilla elevadora), catalogándolo como señal de avería técnica y no como un falso positivo de incendio.

\section{Sistemas de Detección por Aspiración de Alta Sensibilidad (ASD / VESDA)}

Los sistemas de detección por aspiración de muy alta sensibilidad (conocidos ampliamente bajo el estándar VESDA,
representan el escalón tecnológico más avanzado en la protección de ambientes críticos, tales como salas de servidores,
salas limpias de semiconductores, centros de conmutación telefónica y tableros de distribución de media y baja tensión.

El equipamiento consta de una red de cañerías plásticas no propagadoras de llama (CPVC o ABS) que recorre el volumen
protegido, con orificios de muestreo de diámetro calibrado perforados según un cálculo fluido-dinámico previo. Un
aspirador centrífugo de alto rendimiento succiona continuamente muestras de aire de toda la red y las conduce hacia un
módulo central de análisis. Antes de ingresar a la cámara de sensado, el flujo de aire atraviesa un filtro de dos etapas: la
primera retiene partículas de polvo ambiental de gran tamaño, mientras que la segunda inyecta aire ultrafiltrado y seco
para mantener limpios los componentes ópticos internos.

La cámara de detección emplea un láser de estado sólido de alta precisión acoplado a un arreglo de fotodiodos capaces de
realizar el conteo individual de partículas a través de dispersión Mie. Estos equipos alcanzan umbrales de sensibilidad
extraordinarios, regulables en rangos de oscurecimiento de $0.001\percent/\m$ a $20\percent/\m$, lo que equivale a una
sensibilidad varios órdenes de magnitud superior a la de un detector óptico convencional. Esto permite detectar las
primeras moléculas desprendidas durante la pirólisis de aislamientos de cables o sobrecalentamiento de circuitos impresos
horas antes de que se observe humo a simple vista o se produzcan llamas abiertas.

\section{Detectores Térmicos y de Temperatura}

Los detectores térmicos responden a la transferencia de energía calórica liberada por el fuego hacia el sensor por
convección o conducción. Si bien presentan un tiempo de respuesta más lento comparado con los detectores de humo
(requiriendo que el incendio adquiera cierta masa calórica), ofrecen una inmunidad sobresaliente en ambientes con presencia
habitual de vapores de agua, polvos en suspensión, humos de escape de motores o procesos térmicos de fabricación. Se
clasifican en tres modalidades:

\begin{itemize}
    \item Detectores Termostáticos (Temperatura Fija): Se activan cuando la temperatura del ambiente o del elemento
          sensor alcanza un valor umbral predeterminado de fábrica (típicamente $57\Celsius$, $68\Celsius$ o $74\Celsius$).
          Operan mecánicamente mediante discos bimetálicos de resorte de salto repentino, cápsulas con aleaciones metálicas
          eutécticas fusibles que se funden liberando un contacto bajo resorte, o electrónicamente mediante termistores NTC.
    \item Detectores Termovelocimétricos: Responden a la velocidad o tasa temporal de elevación de
          la temperatura, con total independencia del valor inicial. En su concepción electromecánica, constan de una cámara
          de aire sellada por un diafragma flexible de contacto que posee una fuga o escape calibrado. Cuando la temperatura
          se eleva de forma suave (debido a climatización diurna), el aire se expande y escapa por la fuga sin deformar el
          diafragma. Ante un incendio donde la temperatura crece rápidamente a razón de más de $8\Celsius/\mathrm{min}$ a
          $9\Celsius/\mathrm{min}$, la presión interna no alcanza a disiparse por la fuga, venciendo el diafragma y cerrando
          el circuito de alarma. En variantes electrónicas, se emplean dos termistores con constantes de tiempo térmicas
          diferentes procesados por un microcontrolador.
    \item Detectores Térmicos Combinados: Integran en una misma cápsula el sensado termovelocimétrico para detectar fuegos
          de rápida propagación y el sensado de temperatura fija como respaldo ante fuegos de desarrollo térmico gradual.
\end{itemize}

\section{Cable Sensor Térmico Lineal (LHD)}

Para la protección de activos continuos y alargados donde la instalación de detectores puntuales resulta impracticable, se
utiliza el cable sensor térmico lineal. Se presentan dos variantes tecnológicas:

\begin{itemize}
    \item Cable LHD Digital: Formado por dos conductores elásticos de acero trenzados bajo tensión y
          recubiertos por polímero termosensible con punto de fusión calibrado (e.g., $68\Celsius$,
          $88\Celsius$ o $105\Celsius$). El conjunto posee una cubierta exterior de protección mecánica.
          Cuando la temperatura en cualquier punto alcanza el valor de fusión, el polímero funde y los
          conductores entran en cortocircuito. El módulo de control mide la resistencia hasta el punto
          del contacto, determinando la distancia exacta del foco con precisión métrica.
    \item Cable LHD Analógico y de Fibra Óptica: Los analógicos miden la variación resistiva del aislamiento en función
          de la temperatura. Por su parte, los sensores de fibra óptica emplean Reflectometría Óptica
          en el Dominio del Tiempo (OTDR) analizando el esparcimiento inelástico Raman producido por oscilaciones
          moleculares del vidrio, brindando un perfil térmico continuo en tiempo real sobre distancias de hasta
          $10\unit{\kilo\meter}$.
\end{itemize}
Su campo de aplicación primordial comprende las bandejas portacables de potencia en subestaciones, túneles de servicios,
fajas transportadoras de carbón o cereales, transformadores de distribución e interior de tableros eléctricos.

\section{Detectores de Radiación y Llama}

Los detectores de llama responden a la energía electromagnética emitida en las longitudes de onda ópticas (ultravioleta,
visible e infrarroja) por las especies atómicas y moleculares excitadas en el cono de una llama. Operan con una velocidad de
respuesta sumamente rápida, típicamente inferior a $100\ms$, resultando insustituibles en zonas donde un escape de gases o
líquidos inflamables puede provocar una combustión súbita sin generación previa de humo.

\begin{itemize}
    \item Detectores Ultravioleta (UV): Emplean tubos de descarga fotosensibles al rango espectral de $185\,\mathrm{nm}$ a
          $260\,\mathrm{nm}$. La radiación solar que alcanza la superficie terrestre en esta banda es filtrada por la capa de
          ozono atmosférico, lo que otorga al sensor inmunidad a la luz solar directa. Son capaces de detectar llamas de
          hidrocarburos, hidrógeno y metales con gran velocidad. No obstante, son vulnerables a falsas alarmas provocadas por
          arcos eléctricos de soldadura, descargas de rayos y arcos de alta tensión en seccionadores.
    \item Detectores Infrarrojos (IR): Sintonizan fotocélulas piroeléctricas en la banda de absorción y emisión de
          resonancia del dióxido de carbono caliente a $4.3\,\mu\mathrm{m}$, complementada con un filtro pasabanda óptico y
          circuitos que reconocen la frecuencia de parpadeo característica del fuego (frecuencias de modulación de llama
          entre $1\Hz$ y $10\Hz$).
    \item Detectores Multiespectro UV/IR y Triple Infrarrojo (IR3): Para evitar disparos espurios en ambientes industriales
          complejos, los detectores IR3 incorporan tres sensores infrarrojos independientes que monitorean diferentes
          longitudes de onda: la banda activa de emisión de $CO_2$ ($4.3\,\mu\mathrm{m}$) y dos bandas de referencia
          colaterales ($4.0\,\mu\mathrm{m}$ y $4.6\,\mu\mathrm{m}$). Mediante procesamiento digital de señales, el algoritmo
          compara las relaciones de intensidad y las características de modulación temporal, permitiendo discriminar con
          certeza absoluta una llama real frente a fuentes de calor radiante estacionarias, reflectores halógenos o radiación
          solar reflejada en espejos de agua.
\end{itemize}

\section{Detectores de Gases de Combustión}

Constituyen una herramienta complementaria de monitoreo en recintos confinados. Emplean celdas electroquímicas catalíticas
sensibles a concentraciones de monóxido de carbono (CO) en rangos de $0$ a $500\,\text{ppm}$. Dado que la pirólisis de la
mayoría de los materiales orgánicos libera CO antes de que las corrientes de convección térmica tengan la fuerza de elevar
partículas pesadas de humo, estos detectores posibilitan una detección anticipada en conductos de aire y estacionamientos
subterráneos.

En la Tabla \ref{tab.cap5:deteccion} se resume el análisis comparativo de las diferentes tecnologías de detección.

\begin{table}[!ht]
    \centering
    \resizebox{\textwidth}{!}{
    \begin{tabular}{|>{\raggedright\arraybackslash}p{3.5cm}|>{\raggedright\arraybackslash}p{4cm}|>{\raggedright\arraybackslash}p{2.5cm}|>{\raggedright\arraybackslash}p{3cm}|>{\raggedright\arraybackslash}p{4.5cm}|}
        \hline
        \textbf{Tecnología} & \textbf{Principio Físico} & \textbf{Velocidad} & \textbf{Inmunidad a Fallas} & \textbf{Aplicación Recomendada} \\ \hline
        Óptico Fotoeléctrico & Dispersión Tyndall en cámara oscura. & Media & Media (afectado por polvillo y vapor). &
        Oficinas, pasillos, locales comerciales, salas de reunión. \\ \hline
        Iónico & Disminución de corriente iónica por $^{241}\mathrm{Am}$. & Rápida en llama abierta & Baja (sensible a corrientes de aire). &
        Uso en desuso por restricciones ambientales. \\ \hline
        Haz Proyectado (Barrera) & Atenuación infrarroja en trayectoria larga ($10\text{--}100\m$). & Media & Alta con discriminación de bloqueo. &
        Naves industriales de gran altura, depósitos, hangares. \\ \hline
        Aspiración (ASD/VESDA) & Conteo láser Mie de partículas con red de muestreo activa. & Muy rápida (incipiente) & Muy alta con doble filtrado. &
        Centros de datos, salas limpias, racks electrónicos cerrados. \\ \hline
        Térmico Fijo / Velocimétrico & Deformación bimetálica o diferencial termistor. & Lenta a moderada & Excelente (inmune a humo/polvillo). &
        Cocinas, salas de calderas, talleres con humos de proceso. \\ \hline
        Cable LHD & Fusión polimérica o esparcimiento Raman en fibra. & Moderada a lo largo del tendido & Alta ante agresiones químicas y agua. &
        Bandejas de cables, transformadores, túneles de cables. \\ \hline
        Llama (IR3 / UV-IR) & Espectrometría óptica multicanal y parpadeo de llama. & Instantánea ($< 100\ms$) & Muy alta en configuración IR3. &
        Almacenamiento de combustibles, hangares, petroquímica. \\ \hline
    \end{tabular}
    }
    \caption{Matriz comparativa de tecnologías de detección de incendios.}
    \label{tab.cap5:deteccion}
\end{table}
